% ============================================================================
% Pacchetti aggiuntivi compatibili con la classe supsistudent.
% NON ricaricare: babel, geometry, inputenc, fontenc, helvet, setspace,
% graphicx, fancyhdr, titlesec, amsmath, varioref, lastpage, eso-pic,
% changepage  (sono già caricati da supsistudent.cls).
% ============================================================================

% ── Tipografia ─────────────────────────────────────────────────────────────
\usepackage{microtype}

% Tolleranza moderata per spezzare righe difficili senza generare
% troppi "Underfull". Aiuta soprattutto su URL e identificatori lunghi.
\setlength{\emergencystretch}{1.5em}

% ── Matematica & simboli (amsmath già caricato dalla classe) ───────────────
\usepackage{amssymb}
\usepackage{mathtools}
\usepackage{siunitx}
\sisetup{
  detect-all,
  per-mode=symbol,
  output-decimal-marker={.}
}

% ── Tabelle, liste ─────────────────────────────────────────────────────────
\usepackage{booktabs}
\usepackage{tabularx}
\usepackage{longtable}
\usepackage{multirow}
\usepackage{enumitem}
\setlist{topsep=2pt, itemsep=2pt, parsep=2pt}

% ── Colori ─────────────────────────────────────────────────────────────────
\usepackage[table]{xcolor}
\definecolor{accent}{HTML}{1F4E79}
\definecolor{codebg}{HTML}{F4F6F8}
\definecolor{codestring}{HTML}{1A7F37}
\definecolor{codekw}{HTML}{8250DF}
\definecolor{codecomment}{HTML}{6E7781}

% ── Listing di codice ──────────────────────────────────────────────────────
\usepackage{listings}

% YAML (listings non lo supporta nativamente)
\lstdefinelanguage{YAML}{
  keywords={true, false, null},
  keywordstyle=\color{codekw}\bfseries,
  sensitive=true,
  comment=[l]{\#},
  morestring=[b]",
  morestring=[b]',
  literate=
    {:}{{{\color{codekw}{:}}}}{1}
    {-}{{{\color{codekw}{-}}}}{1},
}

% JSON
\lstdefinelanguage{JSON}{
  keywords={true, false, null},
  keywordstyle=\color{codekw}\bfseries,
  sensitive=false,
  morestring=[b]",
  literate=
    {\{}{{{\color{codekw}\{}}}{1}
    {\}}{{{\color{codekw}\}}}}{1}
    {[}{{{\color{codekw}[}}}{1}
    {]}{{{\color{codekw}]}}}{1}
    {:}{{{\color{codekw}{:}}}}{1}
    {,}{{{\color{codekw}{,}}}}{1},
}

\lstset{
  basicstyle=\ttfamily\footnotesize,
  backgroundcolor=\color{codebg},
  frame=single,
  rulecolor=\color{codebg},
  framerule=0pt,
  xleftmargin=8pt,
  xrightmargin=8pt,
  aboveskip=4pt,
  belowskip=4pt,
  showstringspaces=false,
  breaklines=true,
  keywordstyle=\color{codekw}\bfseries,
  commentstyle=\color{codecomment}\itshape,
  stringstyle=\color{codestring},
  language=Python,
  tabsize=4,
}

% ── Figure (graphicx già caricato dalla classe) ────────────────────────────
\graphicspath{{figures/}}
\usepackage{subcaption}
\usepackage{float}
\usepackage{caption}
\captionsetup{font=small, labelfont=bf}

% ── Box pedagogici (Key Insight) ───────────────────────────────────────────
\usepackage[most]{tcolorbox}
\newtcolorbox{insight}[1][]{
  enhanced,
  colback=accent!5!white,
  colframe=accent!75!black,
  fonttitle=\bfseries,
  title={Key Insight},
  rounded corners,
  drop fuzzy shadow=accent!20!gray,
  boxrule=0.5pt,
  #1
}

% ── Riferimenti & link ─────────────────────────────────────────────────────
\usepackage[hidelinks]{hyperref}
\hypersetup{
  pdftitle={Battery GEIS MLOps Platform},
  pdfauthor={Davide Corso, Marco Soldani},
  pdfsubject={Progetto di Semestre — MLOps},
  pdfkeywords={MLOps, EIS, GEIS, Lithium-Ion, Battery, Drift Detection, MLflow}
}
\usepackage[noabbrev,capitalise,italian]{cleveref}

% ── Bibliografia (biblatex + biber) ────────────────────────────────────────
\usepackage{csquotes}
\usepackage[backend=biber,style=numeric,sorting=none,maxbibnames=99,minbibnames=99]{biblatex}

% ── Glossari / acronimi ────────────────────────────────────────────────────
\usepackage[acronym,toc,nonumberlist,nopostdot,nogroupskip]{glossaries}
\setglossarystyle{long}
\makeglossaries

% ── Comandi custom ─────────────────────────────────────────────────────────
% Forzato in math mode: \Omega in Helvetica (font sans default SUPSI) non
% esiste e renderizza come tofu box. In math mode usa cm-math che ce l'ha.
\newcommand{\milliohm}{\ensuremath{\si{\milli\ohm}}}
\newcommand{\codefile}[1]{\texttt{\small#1}}
